% --- Archon physics formalization source begin ---
% archon:physics
% archon:covers PhyXMiniProblems/problem_phyx_mini_0863.lean
% archon:source-report reports/phyx_mini/problem_phyx_mini_0863.source.json

\chapter{Physics problem phyx\_mini\_0863}
\label{ch:PhyXMiniProblems_problem_phyx_mini_0863}

\paragraph{Problem source.}
\#\# Physical scenario

The two halves of the rod in figure are uniformly charged to \textbackslash{}( \textbackslash{}pm Q \textbackslash{}).

\#\# Question

What is the electric potential at the point indicated by the dot?

\#\# Answer choices

- A: A: 1.2V

- B: B: 1.4V

- C: C: 2V

- D: D: 0V

\#\# Figure caption (auxiliary; use the image as primary evidence)

The image features a horizontal rectangular rod divided into two sections. The left section has red plus signs indicating positive charge, and the right section has a teal dashed line indicating neutrality or separation. Above the rectangle, there's a black dot representing a point charge or a point of interest. A vertical arrow labeled "d" points from the rectangle to the dot, suggesting a distance. A horizontal double-headed arrow labeled "L" spans the rectangle's length, indicating its length. The scene is likely a physics illustration related to electric fields or potential.

\#\# Dataset metadata

Electrostatics; Physical Model Grounding Reasoning, Multi-Formula Reasoning

\paragraph{Recorded answer/context.}
D: D: 0V

\paragraph{Figure/image path.}
/root/proposal\_for\_physic/science-mango/phyx\_mini\_run/phyx\_data/test\_image/863.png

\paragraph{Formalization target.}
create a compiling Lean file with sorry bodies at `PhyXMiniProblems/problem\_phyx\_mini\_0863.lean`. The Lean declarations must preserve the physical quantities, dimensions or dimensional roles, figure labels, governing-law hypotheses, and final relation expressed by this problem.
Use Mathlib/Physlib names found through LeanExplore where available. If a physics API is missing, introduce faithful local abstractions rather than scalar placeholder aliases.

\begin{theorem}[Physics formalization target]
\label{thm:physics:phyx_mini_0863:target}
\lean{PhyXMiniProblems.ProblemPhyXMini0863.electricPotentialAtDotOfOppositelyChargedRod}
\uses{def:physics:phyx-mini-0863:phyxminiproblems-problemphyxmini0863-electricpotentialinvolts, def:physics:phyx-mini-0863:phyxminiproblems-problemphyxmini0863-splitrodpotentialsetup, def:physics:phyx-mini-0863:phyxminiproblems-problemphyxmini0863-matchesprimarysplitrodfigure, def:physics:phyx-mini-0863:phyxminiproblems-problemphyxmini0863-hasphysicalsplitrodparameters, def:physics:phyx-mini-0863:phyxminiproblems-problemphyxmini0863-useszeropotentialatspatialinfinity, def:physics:phyx-mini-0863:phyxminiproblems-problemphyxmini0863-satisfiesuniformoppositehalfchargemodel, def:physics:phyx-mini-0863:phyxminiproblems-problemphyxmini0863-satisfiescoulombpotentialsuperpositionlaw, def:physics:phyx-mini-0863:phyxminiproblems-problemphyxmini0863-isuniqueexactpotentialanswer}
The assigned autoformalize agent should translate this physics problem into Lean declarations in the covered file, with theorem and lemma proof bodies written as `by sorry`.
\end{theorem}
\begin{proof}
This is an autoformalization task, not a proof task. Produce statements that can later be proved by the physics prover without weakening the physical model.
\end{proof}

% --- Archon declaration topology begin ---
\section{Declaration topology}
\begin{definition}[linear Charge Density Dimension]
  \label{def:physics:phyx-mini-0863:phyxminiproblems-problemphyxmini0863-linearchargedensitydimension}
  \lean{PhyXMiniProblems.ProblemPhyXMini0863.linearChargeDensityDimension}
  The dimension C L⁻¹ of linear electric-charge density
\end{definition}
\begin{proof}
  This is the declared model definition of linear Charge Density Dimension.
\end{proof}
\begin{definition}[electric Potential Dimension]
  \label{def:physics:phyx-mini-0863:phyxminiproblems-problemphyxmini0863-electricpotentialdimension}
  \lean{PhyXMiniProblems.ProblemPhyXMini0863.electricPotentialDimension}
  The dimension M L² T⁻² C⁻¹ of electric potential
\end{definition}
\begin{proof}
  This is the declared model definition of electric Potential Dimension.
\end{proof}
\begin{definition}[Length Quantity]
  \label{def:physics:phyx-mini-0863:phyxminiproblems-problemphyxmini0863-lengthquantity}
  \lean{PhyXMiniProblems.ProblemPhyXMini0863.LengthQuantity}
  A nonnegative, unit-independent physical length
\end{definition}
\begin{proof}
  This is the declared model definition of Length Quantity.
\end{proof}
\begin{definition}[Charge Magnitude Quantity]
  \label{def:physics:phyx-mini-0863:phyxminiproblems-problemphyxmini0863-chargemagnitudequantity}
  \lean{PhyXMiniProblems.ProblemPhyXMini0863.ChargeMagnitudeQuantity}
  A nonnegative, unit-independent magnitude of electric charge
\end{definition}
\begin{proof}
  This is the declared model definition of Charge Magnitude Quantity.
\end{proof}
\begin{definition}[Linear Charge Density Quantity]
  \label{def:physics:phyx-mini-0863:phyxminiproblems-problemphyxmini0863-linearchargedensityquantity}
  \lean{PhyXMiniProblems.ProblemPhyXMini0863.LinearChargeDensityQuantity}
  \uses{def:physics:phyx-mini-0863:phyxminiproblems-problemphyxmini0863-linearchargedensitydimension}
  A signed, unit-independent physical linear charge density
\end{definition}
\begin{proof}
  This is the declared model definition of Linear Charge Density Quantity.
\end{proof}
\begin{definition}[Electric Potential Quantity]
  \label{def:physics:phyx-mini-0863:phyxminiproblems-problemphyxmini0863-electricpotentialquantity}
  \lean{PhyXMiniProblems.ProblemPhyXMini0863.ElectricPotentialQuantity}
  \uses{def:physics:phyx-mini-0863:phyxminiproblems-problemphyxmini0863-electricpotentialdimension}
  A signed, unit-independent physical electric potential
\end{definition}
\begin{proof}
  This is the declared model definition of Electric Potential Quantity.
\end{proof}
\begin{definition}[length In Meters]
  \label{def:physics:phyx-mini-0863:phyxminiproblems-problemphyxmini0863-lengthinmeters}
  \lean{PhyXMiniProblems.ProblemPhyXMini0863.lengthInMeters}
  \uses{def:physics:phyx-mini-0863:phyxminiproblems-problemphyxmini0863-lengthquantity}
  Read a physical length in coherent-SI metres
\end{definition}
\begin{proof}
  This is the declared model definition of length In Meters.
\end{proof}
\begin{definition}[charge Magnitude In Coulombs]
  \label{def:physics:phyx-mini-0863:phyxminiproblems-problemphyxmini0863-chargemagnitudeincoulombs}
  \lean{PhyXMiniProblems.ProblemPhyXMini0863.chargeMagnitudeInCoulombs}
  \uses{def:physics:phyx-mini-0863:phyxminiproblems-problemphyxmini0863-chargemagnitudequantity}
  Read a physical charge magnitude in coherent-SI coulombs
\end{definition}
\begin{proof}
  This is the declared model definition of charge Magnitude In Coulombs.
\end{proof}
\begin{definition}[linear Charge Density In Coulombs Per Meter]
  \label{def:physics:phyx-mini-0863:phyxminiproblems-problemphyxmini0863-linearchargedensityincoulombspermeter}
  \lean{PhyXMiniProblems.ProblemPhyXMini0863.linearChargeDensityInCoulombsPerMeter}
  \uses{def:physics:phyx-mini-0863:phyxminiproblems-problemphyxmini0863-linearchargedensityquantity}
  Read a signed line-charge density in coherent-SI coulombs per metre
\end{definition}
\begin{proof}
  This is the declared model definition of linear Charge Density In Coulombs Per Meter.
\end{proof}
\begin{definition}[electric Potential In Volts]
  \label{def:physics:phyx-mini-0863:phyxminiproblems-problemphyxmini0863-electricpotentialinvolts}
  \lean{PhyXMiniProblems.ProblemPhyXMini0863.electricPotentialInVolts}
  \uses{def:physics:phyx-mini-0863:phyxminiproblems-problemphyxmini0863-electricpotentialquantity}
  Read a signed electric potential in coherent-SI volts
\end{definition}
\begin{proof}
  This is the declared model definition of electric Potential In Volts.
\end{proof}
\begin{definition}[Uniform Oppositely Charged Rod]
  \label{def:physics:phyx-mini-0863:phyxminiproblems-problemphyxmini0863-uniformoppositelychargedrod}
  \lean{PhyXMiniProblems.ProblemPhyXMini0863.UniformOppositelyChargedRod}
  \uses{def:physics:phyx-mini-0863:phyxminiproblems-problemphyxmini0863-lengthquantity, def:physics:phyx-mini-0863:phyxminiproblems-problemphyxmini0863-chargemagnitudequantity, def:physics:phyx-mini-0863:phyxminiproblems-problemphyxmini0863-linearchargedensityquantity}
  The split rod and its independent charge observables. The argument of linearChargeDensityAt is an axial coordinate read in SI metres
\end{definition}
\begin{proof}
  This is the declared model definition of Uniform Oppositely Charged Rod.
\end{proof}
\begin{definition}[Figure Point]
  \label{def:physics:phyx-mini-0863:phyxminiproblems-problemphyxmini0863-figurepoint}
  \lean{PhyXMiniProblems.ProblemPhyXMini0863.FigurePoint}
  Distinguished points shown in, or determined by, the primary image
\end{definition}
\begin{proof}
  This is the declared model definition of Figure Point.
\end{proof}
\begin{definition}[Figure Feature]
  \label{def:physics:phyx-mini-0863:phyxminiproblems-problemphyxmini0863-figurefeature}
  \lean{PhyXMiniProblems.ProblemPhyXMini0863.FigureFeature}
  Qualitative features visible in the supplied bitmap 863.png
\end{definition}
\begin{proof}
  This is the declared model definition of Figure Feature.
\end{proof}
\begin{definition}[Split Charged Rod Figure]
  \label{def:physics:phyx-mini-0863:phyxminiproblems-problemphyxmini0863-splitchargedrodfigure}
  \lean{PhyXMiniProblems.ProblemPhyXMini0863.SplitChargedRodFigure}
  \uses{def:physics:phyx-mini-0863:phyxminiproblems-problemphyxmini0863-lengthquantity, def:physics:phyx-mini-0863:phyxminiproblems-problemphyxmini0863-figurepoint, def:physics:phyx-mini-0863:phyxminiproblems-problemphyxmini0863-figurefeature}
  The figure is represented in an SI-metre coordinate chart. Its physical labels L and d remain dimensionful quantities rather than scalar aliases
\end{definition}
\begin{proof}
  This is the declared model definition of Split Charged Rod Figure.
\end{proof}
\begin{definition}[Electric Potential Reference]
  \label{def:physics:phyx-mini-0863:phyxminiproblems-problemphyxmini0863-electricpotentialreference}
  \lean{PhyXMiniProblems.ProblemPhyXMini0863.ElectricPotentialReference}
  The reference convention for the scalar electric potential
\end{definition}
\begin{proof}
  This is the declared model definition of Electric Potential Reference.
\end{proof}
\begin{definition}[Split Rod Potential Setup]
  \label{def:physics:phyx-mini-0863:phyxminiproblems-problemphyxmini0863-splitrodpotentialsetup}
  \lean{PhyXMiniProblems.ProblemPhyXMini0863.SplitRodPotentialSetup}
  \uses{def:physics:phyx-mini-0863:phyxminiproblems-problemphyxmini0863-lengthquantity, def:physics:phyx-mini-0863:phyxminiproblems-problemphyxmini0863-electricpotentialquantity, def:physics:phyx-mini-0863:phyxminiproblems-problemphyxmini0863-uniformoppositelychargedrod, def:physics:phyx-mini-0863:phyxminiproblems-problemphyxmini0863-splitchargedrodfigure, def:physics:phyx-mini-0863:phyxminiproblems-problemphyxmini0863-electricpotentialreference}
  The independent physical setup. In particular, the potential at the dot is stored as an observable; it is not defined to be any displayed answer value
\end{definition}
\begin{proof}
  This is the declared model definition of Split Rod Potential Setup.
\end{proof}
\begin{definition}[Matches Primary Split Rod Figure]
  \label{def:physics:phyx-mini-0863:phyxminiproblems-problemphyxmini0863-matchesprimarysplitrodfigure}
  \lean{PhyXMiniProblems.ProblemPhyXMini0863.MatchesPrimarySplitRodFigure}
  \uses{def:physics:phyx-mini-0863:phyxminiproblems-problemphyxmini0863-lengthinmeters, def:physics:phyx-mini-0863:phyxminiproblems-problemphyxmini0863-splitrodpotentialsetup}
  Transcription of the primary image: the division is the coordinate origin, the rod endpoints are at ±L/2, and the dot is at (0,d). This predicate contains no potential value or answer choice
\end{definition}
\begin{proof}
  This is the declared model definition of Matches Primary Split Rod Figure.
\end{proof}
\begin{definition}[Has Physical Split Rod Parameters]
  \label{def:physics:phyx-mini-0863:phyxminiproblems-problemphyxmini0863-hasphysicalsplitrodparameters}
  \lean{PhyXMiniProblems.ProblemPhyXMini0863.HasPhysicalSplitRodParameters}
  \uses{def:physics:phyx-mini-0863:phyxminiproblems-problemphyxmini0863-lengthinmeters, def:physics:phyx-mini-0863:phyxminiproblems-problemphyxmini0863-chargemagnitudeincoulombs, def:physics:phyx-mini-0863:phyxminiproblems-problemphyxmini0863-splitrodpotentialsetup}
  Positivity assumptions selecting the nondegenerate depicted setup
\end{definition}
\begin{proof}
  This is the declared model definition of Has Physical Split Rod Parameters.
\end{proof}
\begin{definition}[Uses Zero Potential At Spatial Infinity]
  \label{def:physics:phyx-mini-0863:phyxminiproblems-problemphyxmini0863-useszeropotentialatspatialinfinity}
  \lean{PhyXMiniProblems.ProblemPhyXMini0863.UsesZeroPotentialAtSpatialInfinity}
  \uses{def:physics:phyx-mini-0863:phyxminiproblems-problemphyxmini0863-splitrodpotentialsetup}
  The conventional electrostatic choice of zero potential at infinity
\end{definition}
\begin{proof}
  This is the declared model definition of Uses Zero Potential At Spatial Infinity.
\end{proof}
\begin{definition}[rod Line Charge Density In Coulombs Per Meter]
  \label{def:physics:phyx-mini-0863:phyxminiproblems-problemphyxmini0863-rodlinechargedensityincoulombspermeter}
  \lean{PhyXMiniProblems.ProblemPhyXMini0863.rodLineChargeDensityInCoulombsPerMeter}
  \uses{def:physics:phyx-mini-0863:phyxminiproblems-problemphyxmini0863-linearchargedensityincoulombspermeter, def:physics:phyx-mini-0863:phyxminiproblems-problemphyxmini0863-splitrodpotentialsetup}
  SI line-density readout at the axial coordinate xMeters
\end{definition}
\begin{proof}
  This is the declared model definition of rod Line Charge Density In Coulombs Per Meter.
\end{proof}
\begin{definition}[Satisfies Uniform Opposite Half Charge Model]
  \label{def:physics:phyx-mini-0863:phyxminiproblems-problemphyxmini0863-satisfiesuniformoppositehalfchargemodel}
  \lean{PhyXMiniProblems.ProblemPhyXMini0863.SatisfiesUniformOppositeHalfChargeModel}
  \uses{def:physics:phyx-mini-0863:phyxminiproblems-problemphyxmini0863-lengthinmeters, def:physics:phyx-mini-0863:phyxminiproblems-problemphyxmini0863-chargemagnitudeincoulombs, def:physics:phyx-mini-0863:phyxminiproblems-problemphyxmini0863-splitrodpotentialsetup, def:physics:phyx-mini-0863:phyxminiproblems-problemphyxmini0863-rodlinechargedensityincoulombspermeter}
  Each half has length L/2 and total charge of magnitude Q, so its uniform density has magnitude 2Q/L. The endpoint convention at the single division point is immaterial to the line integrals. The two integral fields explicitly record that the half-rod totals are +Q and -Q
\end{definition}
\begin{proof}
  This is the declared model definition of Satisfies Uniform Opposite Half Charge Model.
\end{proof}
\begin{definition}[source To Dot Distance In Meters]
  \label{def:physics:phyx-mini-0863:phyxminiproblems-problemphyxmini0863-sourcetodotdistanceinmeters}
  \lean{PhyXMiniProblems.ProblemPhyXMini0863.sourceToDotDistanceInMeters}
  \uses{def:physics:phyx-mini-0863:phyxminiproblems-problemphyxmini0863-lengthinmeters, def:physics:phyx-mini-0863:phyxminiproblems-problemphyxmini0863-splitrodpotentialsetup}
  Euclidean source-to-dot distance for a source element at axial coordinate xMeters, using the figure coordinates (x,0) and (0,d)
\end{definition}
\begin{proof}
  This is the declared model definition of source To Dot Distance In Meters.
\end{proof}
\begin{definition}[potential Contribution Density In Volts Per Meter]
  \label{def:physics:phyx-mini-0863:phyxminiproblems-problemphyxmini0863-potentialcontributiondensityinvoltspermeter}
  \lean{PhyXMiniProblems.ProblemPhyXMini0863.potentialContributionDensityInVoltsPerMeter}
  \uses{def:physics:phyx-mini-0863:phyxminiproblems-problemphyxmini0863-splitrodpotentialsetup, def:physics:phyx-mini-0863:phyxminiproblems-problemphyxmini0863-rodlinechargedensityincoulombspermeter, def:physics:phyx-mini-0863:phyxminiproblems-problemphyxmini0863-sourcetodotdistanceinmeters}
  The coherent-SI scalar-potential contribution density k λ/r, in volts per metre, for a rod element at xMeters. Physlib supplies Coulomb's constant as Electromagnetism.EMSystem.coulombConstant
\end{definition}
\begin{proof}
  This is the declared model definition of potential Contribution Density In Volts Per Meter.
\end{proof}
\begin{definition}[Satisfies Coulomb Potential Superposition Law]
  \label{def:physics:phyx-mini-0863:phyxminiproblems-problemphyxmini0863-satisfiescoulombpotentialsuperpositionlaw}
  \lean{PhyXMiniProblems.ProblemPhyXMini0863.SatisfiesCoulombPotentialSuperpositionLaw}
  \uses{def:physics:phyx-mini-0863:phyxminiproblems-problemphyxmini0863-lengthinmeters, def:physics:phyx-mini-0863:phyxminiproblems-problemphyxmini0863-electricpotentialinvolts, def:physics:phyx-mini-0863:phyxminiproblems-problemphyxmini0863-splitrodpotentialsetup, def:physics:phyx-mini-0863:phyxminiproblems-problemphyxmini0863-potentialcontributiondensityinvoltspermeter}
  Coulomb superposition for the scalar potential, with the rod supported on [-L/2,L/2]. This is a general source integral and does not state its value
\end{definition}
\begin{proof}
  This is the declared model definition of Satisfies Coulomb Potential Superposition Law.
\end{proof}
\begin{lemma}[source To Dot Distance reflection]
  \label{lem:physics:phyx-mini-0863:phyxminiproblems-problemphyxmini0863-sourcetodotdistance-reflection}
  \lean{PhyXMiniProblems.ProblemPhyXMini0863.sourceToDotDistance_reflection}
  \uses{def:physics:phyx-mini-0863:phyxminiproblems-problemphyxmini0863-splitrodpotentialsetup, def:physics:phyx-mini-0863:phyxminiproblems-problemphyxmini0863-sourcetodotdistanceinmeters}
  Reflection across the division leaves the source-to-dot distance fixed
\end{lemma}
\begin{proof}
  The conclusion follows from the governing relations and figure readouts named in the statement of source To Dot Distance reflection.
\end{proof}
\begin{lemma}[paired Potential Contributions cancel]
  \label{lem:physics:phyx-mini-0863:phyxminiproblems-problemphyxmini0863-pairedpotentialcontributions-cancel}
  \lean{PhyXMiniProblems.ProblemPhyXMini0863.pairedPotentialContributions_cancel}
  \uses{def:physics:phyx-mini-0863:phyxminiproblems-problemphyxmini0863-lengthinmeters, def:physics:phyx-mini-0863:phyxminiproblems-problemphyxmini0863-splitrodpotentialsetup, def:physics:phyx-mini-0863:phyxminiproblems-problemphyxmini0863-satisfiesuniformoppositehalfchargemodel, def:physics:phyx-mini-0863:phyxminiproblems-problemphyxmini0863-potentialcontributiondensityinvoltspermeter}
  Opposite source elements have cancelling potential contributions
\end{lemma}
\begin{proof}
  The conclusion follows from the governing relations and figure readouts named in the statement of paired Potential Contributions cancel.
\end{proof}
\begin{definition}[Answer Choice]
  \label{def:physics:phyx-mini-0863:phyxminiproblems-problemphyxmini0863-answerchoice}
  \lean{PhyXMiniProblems.ProblemPhyXMini0863.AnswerChoice}
  Labels of the four voltage answers printed in the question
\end{definition}
\begin{proof}
  This is the declared model definition of Answer Choice.
\end{proof}
\begin{definition}[potential In Volts]
  \label{def:physics:phyx-mini-0863:phyxminiproblems-problemphyxmini0863-answerchoice-potentialinvolts}
  \lean{PhyXMiniProblems.ProblemPhyXMini0863.AnswerChoice.potentialInVolts}
  \uses{def:physics:phyx-mini-0863:phyxminiproblems-problemphyxmini0863-answerchoice}
  Electric potential printed beside each choice, in volts
\end{definition}
\begin{proof}
  This is the declared model definition of potential In Volts.
\end{proof}
\begin{definition}[Is Unique Exact Potential Answer]
  \label{def:physics:phyx-mini-0863:phyxminiproblems-problemphyxmini0863-isuniqueexactpotentialanswer}
  \lean{PhyXMiniProblems.ProblemPhyXMini0863.IsUniqueExactPotentialAnswer}
  \uses{def:physics:phyx-mini-0863:phyxminiproblems-problemphyxmini0863-electricpotentialinvolts, def:physics:phyx-mini-0863:phyxminiproblems-problemphyxmini0863-splitrodpotentialsetup, def:physics:phyx-mini-0863:phyxminiproblems-problemphyxmini0863-answerchoice}
  A choice is the unique printed value equal to the physical potential
\end{definition}
\begin{proof}
  This is the declared model definition of Is Unique Exact Potential Answer.
\end{proof}
\begin{definition}[recorded Dataset Answer]
  \label{def:physics:phyx-mini-0863:phyxminiproblems-problemphyxmini0863-recordeddatasetanswer}
  \lean{PhyXMiniProblems.ProblemPhyXMini0863.recordedDatasetAnswer}
  \uses{def:physics:phyx-mini-0863:phyxminiproblems-problemphyxmini0863-answerchoice}
  The source dataset's recorded answer, retained only as unused metadata
\end{definition}
\begin{proof}
  This is the declared model definition of recorded Dataset Answer.
\end{proof}
% --- Archon declaration topology end ---
% --- Archon physics formalization source end ---
